% Figure: the Rayleigh quotient of a 2x2 symmetric matrix as a function
% of the input angle, oscillating between lambda_min and lambda_max and
% touching each extreme exactly at the corresponding eigenvector.
\begin{figure}[htbp]
\centering
\begin{tikzpicture}
\begin{axis}[
  width=13cm,
  height=8cm,
  xlabel={input angle $\theta$ (degrees)},
  ylabel={Rayleigh quotient $R(\vect{x}(\theta))$},
  xmin=0, xmax=180,
  ymin=0.5, ymax=5.5,
  xtick={0,45,90,135,180},
  ytick={1,2,3,4,5},
  tick label style={font=\footnotesize},
  label style={font=\small},
  legend style={font=\footnotesize, at={(0.5,-0.18)}, anchor=north, legend columns=2},
  title={Rayleigh quotient of $\protect\begin{psmallmatrix}4 \& 1 \\ 1 \& 4\protect\end{psmallmatrix}$ against input direction},
  title style={font=\bfseries},
  ymajorgrids=true,
  grid style={gray!25},
]
% R(theta) = 4 + sin(2 theta) for A = [[4,1],[1,4]] with unit vector
% (cos t, sin t): x'Ax = 4 + 2 cos t sin t = 4 + sin(2t).
\addplot[engblue, very thick, domain=0:180, samples=181] {4 + sin(2*x)};
\addlegendentry{$R(\vect{x}(\theta))$}
% lambda_max = 5.
\addplot[engred, thick, dashed, domain=0:180] {5};
\addlegendentry{$\lambda_{\max} = 5$}
% lambda_min = 3.
\addplot[engamber, thick, dashed, domain=0:180] {3};
\addlegendentry{$\lambda_{\min} = 3$}
% Touch points: max at 45 deg (eigenvector (1,1)), min at 135 deg.
\addplot[engred, only marks, mark=*, mark size=2pt] coordinates {(45,5)};
\addplot[engamber, only marks, mark=*, mark size=2pt] coordinates {(135,3)};
\node[font=\footnotesize, engred, anchor=south] at (axis cs:45,5.05)
  {$\vect{v}_{\max}=(1,1)/\sqrt{2}$};
\node[font=\footnotesize, engamber, anchor=north] at (axis cs:135,2.95)
  {$\vect{v}_{\min}=(1,-1)/\sqrt{2}$};
\end{axis}
\end{tikzpicture}
\caption[Rayleigh quotient between the extreme eigenvalues]{The
Rayleigh quotient $R(\vect{x}) = \vect{x}\transpose\mat{A}\vect{x}/
\vect{x}\transpose\vect{x}$ of the symmetric matrix
$\protect\begin{psmallmatrix}4 \& 1\\ 1 \& 4\protect\end{psmallmatrix}$,
plotted as the input unit vector sweeps through all directions. The
quotient stays trapped between the smallest eigenvalue $\lambda_{\min}
= 3$ (amber) and the largest $\lambda_{\max} = 5$ (red), touching each
extreme exactly at the corresponding eigenvector and nowhere else. The
curve is flat, to first order, at each touch point, which is the
stationarity that gives the Rayleigh quotient its quadratic accuracy:
an input a few degrees off the eigenvector still returns an eigenvalue
estimate good to the square of that angular error. This variational
trapping is the engine behind the power method's eigenvalue estimate
and behind the Courant-Fischer proof of the spectral theorem.}
\label{fig:vol02:ch10:rayleigh-quotient}
\end{figure}
